\chapter*{ABSTRACT}

Effective triage is critical for matching hospital resources to patient acuity. In low- and middle-income countries, including Ghana, hospitals frequently suffer from flow management failures characterized by overcrowding, delayed initial assessments, and acuity discontinuity. To address these operational bottlenecks, this project presents a hospital chatbot powered by a fine-tuned Biomedical Bidirectional Encoder Representations from Transformers (BioBERT) model. The system estimates Emergency Severity Index (ESI)--style urgency from unstructured chief complaints and applies a proposed deterministic map \(f_{\mathrm{SATS}}\) to generate SATS-aligned colour and pathway recommendations.

The underlying mathematical architecture uses tokenisation, continuous vector embeddings, and multi-head self-attention to generate contextualized text summaries, which are passed through a Softmax classification head to predict a categorical distribution over five mutually exclusive ESI-style urgency levels. The heuristic \(f_{\mathrm{SATS}}\) then translates these predictions into triage routing colours (Red, Orange, Yellow, or Green); it is an operational crosswalk, not a claim of clinical ESI--SATS equivalence. To support clinical safety, the pipeline features a mathematical relevance gate to filter non-medical inputs prior to classification. Furthermore, the model acts as a text-first precursor to the standard Triage Early Warning Score (TEWS), supplying an initial pathway card \(P(C)\), detailing destination and maximum wait time, before physical vitals can be obtained.

Primary reported metrics are derived from external evaluation on the Medical Information Mart for Intensive Care IV Emergency Department module (MIMIC-IV-ED) and the National Hospital Ambulatory Medical Care Survey (NHAMCS) 2018--22, ensuring evaluation against real-world clinical phrasing rather than synthetic holdouts. On the NHAMCS evaluation set, the baseline BioBERT classifier achieved a five-class accuracy of 47.60\% and a 4-colour mapping accuracy of 49.95\%, with a median inference latency of approximately 94\,ms. To address severe class imbalance and study the safety trade-off between critical-case sensitivity and overall precision, stratified oversampling (text augmentation plus random minority oversampling) was applied during training. This raised Level~1 (Red) recall from 8.00\% to 21.71\% under external domain shift, though at a corresponding cost to overall accuracy. The system functions strictly as a rapid, nurse-overridable decision-support tool, expediting initial routing while preserving human clinical judgment and objective vital-sign scoring.
